../cictikz/src/cictikz/data/tex/cictikz_preamble.tex

% Common source, small signal. The source is grounded, so v_gs = v_i, and
% the output rail carries g_m v_gs and r_ds to ground.

\begin{circuitikz}[american, thick, transform shape, circuit ee IEC]

  ../cictikz/src/cictikz/data/tex/cictikz_lib.tex

  \def\ytop{2.6}
  \def\xgm{3.8}
  \def\xrds{5.6}
  \def\xterm{6.8}

  % input source at the far left
  \draw (0,\ytop) to [V, l_=$v_i$] (0,0);
  \draw (0,\ytop) -- (1.3,\ytop);

  % floating gate marker with v_gs annotation
  \draw (1.39,\ytop) circle (0.09);
  \node[anchor=south] at (1.2,{\ytop+0.1}) {G};
  \node[anchor=north] at (1.4,{\ytop-0.25}) {$+$};
  \node[anchor=west, color=blue] at (0.9,1.3) {$v_{gs}$};
  \node[anchor=south] at (1.4,0.25) {$-$};

  % transconductance and output resistance
  \draw (\xgm,\ytop) to [cI, l_=$g_m v_{gs}$] (\xgm,0);
  \draw (\xrds,\ytop) to [R, l_=$r_{ds}$] (\xrds,0);

  % output rail (drain)
  \draw (\xgm,\ytop) -- (\xterm,\ytop);
  \fill (\xrds,\ytop) circle (0.075);
  \draw (\xterm,\ytop) \portOut{$v_o$};

  % ground rail
  \draw (0,0) -- (\xterm,0);
  \fill (\xgm,0) circle (0.075);
  \fill (\xrds,0) circle (0.075);
  \draw (2.2,0) \vground;

\end{circuitikz}
\end{document}
